% language=us

\environment context-2022-style

\startdocument
  [title={Stepping up \MetaPost},
   banner={convenience and performance},
   location={context\enspace {\bf 2022}\enspace meeting}]

\starttitle[title=\METAPOST\ version 2 versus 3]

\startitemize
    \startitem
        There is less need to adapt \METAPOST\ than \TEX\ to current
        developments, most needed additional can be dealt with in the backend via
        the generic pre- and postscript properties. These properties have been
        added to some more data structures.
    \stopitem
    \startitem
        Some features of \METAPOST\ are not needed or limited so these are
        removed: fonts, backends. Only three number models are supported by
        default: legacy scaled, double, and decimal. We default to double, which
        also performs best.
    \stopitem
    \startitem
        Most significate are the extensions in the \LUA\ interface where we
        got powerful and fast scanners. This permits us to extend the
        language.
    \stopitem
    \startitem
        However, we also also added some primitive features, simply because it
        removed some persistent bottlenecks. The internals didn't really change
        much conceptually. There are numerous optimizations, so it runs faster
        anyway.
    \stopitem
    \startitem
        For Mikael and me this was a nice distraction from dealing with the messy
        math fonts.
    \stopitem
\stopitemize

\stoptitle

\starttitle[title=Intersection]

Getting all intersections of two paths is a bit difficult. It also involves
repetitive restart and cutting before and after and hoping for the best. This is why
we now have

\starttyping
<path> intersectiontimeslist <path>
\stoptyping

This operator returns a path made up of points that reflect the times on the
paths.

\starttyping
intersections-001.tex
\stoptyping

\stoptitle

\starttitle[title=Loops]

\startitemize
    \startitem
        A path is a linked list which means that getting to a point (or between
        points) is a linear lookup.
    \stopitem
    \startitem
        We now have a double linked list and support for progressive loops over
        paths.
    \stopitem
    \startitem
        It is not pretty but works ok for what we need and often the dirty
        details are hidden in macros.
    \stopitem
    \startitem
        New primitives are: \type {deltapoint}, \type {deltaprecontrol}, \type
        {deltapostcontrol}, \type {deltadirection}, \type {pathpoint}, \type
        {pathprecontrol}, \type {pathpostcontrol}, \type {pathdirection}.
    \stopitem
\stopitemize

\starttyping
intersections-001.tex
\stoptyping

\stoptitle

\starttitle[title=Arcs]

The agitate article by Fabrice Larribe was a good reason to see if we could bring
down many hours of rendering to minutes. For that we came up with \type
{arcpoint}, \type {arcpointlist} and \type {subarclength}.

\startbuffer
path p ; p := arcpointlist 100 of (fullcircle xyscaled (10cm,2cm)) ;
draw image ( for i within p :
   draw pathpoint withpen pencircle scaled 2mm ;
endfor ; ) withcolor "extracolor" ;
\stopbuffer

\typebuffer

This gives equally spread points as path:

\startlinecorrection
\processMPbuffer
\stoplinecorrection

\starttyping
agitate-002.tex
ontarget-metapost.tex
\stoptyping

\stoptitle

\starttitle[title=Drawing features]

Instead of grouping with interims for line properties we can also use:

\starttyping
<path> withlinecap <numeric>
<path> withlinejoin <numeric>
<path> withlinemiterlimit <numeric>
\stoptyping

This gives less overhead and is more convenient in generated \METAPOST\ code.

\stoptitle

\starttitle[title=Paths]
The macro \type {center} has two primitive companions:

\starttyping
centerof <path>
centerofmass <path> (Mikael can explain that one)
\stoptyping

These avoid multiple calculations of corners of the bounding box we have:

\starttyping
corners <path|picture>
boundingpath <pen> of <path>
xrange <path|picture>
yrange <path|picture>
\stoptyping

% \startMPcode
% draw(centerofmass (fullcircle xysized(3cm,1cm))) withpen pencircle scaled 3mm withcolor white ;
% draw(center origin withpen pencircle) scaled 3mm withcolor white ;
% draw(centerofmass origin withpen pencircle) scaled 3mm withcolor white ;
% draw(corners (fullcircle xysized(3cm,1cm))) withcolor white ;
% draw(boundingpath pencircle scaled 3mm of (fullcircle xysized(3cm,1cm))) withcolor white;
% \stopMPcode

For generated code or less complex macros we have:

\starttyping
nocycle
\stoptyping

and we have \type {&&} for connecting no|-|continuous paths (as in fonts)
but that is still experimental.

\stoptitle

\starttitle[title=Pens]

Pens can be tricky: circular pens are native \POSTSCRIPT\ (and \PDF) but when
transformed are handled by double paths. Transparency can uncover all kind of
side effects that are normally hidden. The envelopes made by pens are
approximations that assume that these inaccuracies end up inside the filled
areas. Pens get convexed.

\startbuffer
fill convexed (fullcircle xysized (10cm,1cm))
    withcolor "extracolor" withpen pencircle scaled 1mm ;
\stopbuffer

\typebuffer

\startlinecorrection
\processMPbuffer
\stoplinecorrection

The problem is that once we figure this out and apply it low level it's easy to
forget all the details. This is not everyday code and it's hard to make examples.

A good way out of inaccuracies is to rotate over eps which avoids some zero
conditions. The \type {envelope} primitive has been improved and for special
tracing \type {nep} and \type {makenep} are introduced to get around some
heuristics.

% \startMPcode
% pen p ; p := makepen(starring(-1/3) rotated eps) ;
% % nep p ; p := makenep(starring(-1/3) rotated eps) ;
% draw
%     envelope (p scaled 8mm) of (fullcircle xyscaled (.8TextWidth,2cm))
%     withpen pencircle scaled 2mm
%     withcolor darkgreen
%     withtransparency (1,.5) ;
% draw envelope (p scaled 8mm) of reverse (fullcircle xyscaled (.8TextWidth,2cm))
%     withpen pencircle scaled 2mm
%     withcolor darkblue
%     withtransparency (1,.5) ;
% \stopMPcode

\starttyping
ontarget-envelopes.tex
\stoptyping

\stoptitle

\starttitle[title=Lua]

There are plenty of new|/|improved scanners and injectors so that we can build
interfaces that are very efficient. Given the performance of \METAPOST\ we can
extend in \LUA. In \CONTEXT\ most interfaces between \TEX\ and \METAPOST\ are
upgraded.

\stoptitle

\starttitle[title=Etcetera]

We just mention \type {overload}, \type {stacking}, \type {uncontrolled} and
\type {setgroup} and there are some performance related settings. And we're
adding some more in the coming years, like additional support for envelopes.

{\em Mikael Sundqvist made some beautiful, more detailed, overviews of features!}

\stoptitle

\stoptext
